% Independently audited in full by root, B, and D before integration.
\begin{lemma}[Proved: formal lift group for a punctured surface cover]
\label{D:formal-torus-lift}
Let $k$ be algebraically closed of characteristic $p>0$, let
$R=\bigoplus_{d\geq0}R_d$ be a finitely generated normal graded
$k$-domain of dimension two with $R_0=k$, and let $A$ be its
completion at $R_+$.  Let $B/A$ be the normal extension belonging
to a connected finite Galois cover of the punctured spectrum of
$A$, with group $G$.  The grading action of $\mathbb G_m$ on
$A$ lifts to a continuous action on $B$ after a finite \'etale
isogeny $\mathbb G_m\to\mathbb G_m$.  This lifted torus
centralizes the deck group $G$.
\end{lemma}
\begin{proof}
We first construct the lift group on all test algebras; this
avoids assuming that a pointwise Isom functor is of finite type.
Choose homogeneous parameters $x,y\in R_+$ such that $R$ is
finite over $k[x,y]$, and set $C=k[[x,y]]$.  The rings $A$ and
$B$ are complete normal local rings of dimension two, with
residue field $k$.  They are Cohen--Macaulay and finite free over
$C$, by miracle flatness
\cite[tags 0B3D and 00R4]{C_StacksFormalLift}.
For a $k$-algebra $T$ put
\[
 A_T=A\otimes_C T[[x,y]],\qquad
 B_T=B\otimes_C T[[x,y]].
\]
These are also the completed scalar extensions
$\varprojlim_r(A/\mathfrak m_A^r)\otimes_kT$ and
$\varprojlim_r(B/\mathfrak m_B^r)\otimes_kT$:
the parameter-adic and maximal-ideal-adic topologies are cofinal,
and finite freeness over $C$ identifies the parameter-adic limits.

For each $r$, consider pairs consisting of $c\in\mathbb G_m(T)$
and an augmented $T$-algebra automorphism of
$(B/\mathfrak m_B^r)\otimes_kT$ whose restriction to $A$
is the grading action of $c$, followed by the quotient map.
This is an affine group scheme $E_r$ of finite type over $k$.
Indeed, the automorphisms are matrices preserving a finite
multiplication table, the unit, and the augmentation, with
invertible determinant.  Compatibility with $A$ need only be
imposed on finitely many homogeneous algebra generators of $R$;
these are polynomial equations in the matrix entries and
$c,c^{-1}$.  Augmented automorphisms preserve every power of the
augmentation ideal.  Consequently the truncation maps define
an inverse system, and
\[
 E=\varprojlim_r E_r
\]
is an affine group scheme.  Its $T$-points are exactly the pairs
$(c,\phi)$ with $\phi$ a continuous augmentation-preserving
automorphism of $B_T$ lifting $c$ on $A_T$.  Write
$\rho:E\to\mathbb G_m$ for the projection.

We claim that $\ker\rho$ is the constant finite group scheme
$G$, including on nonreduced test algebras $T$.  Let
\[
 U_T=\operatorname{Spec}(A_T)\setminus V(x,y),\qquad
 V_T=\operatorname{Spec}(B_T)\setminus V(x,y).
\]
The map $V_T\to U_T$ is the base change of the given $G$-torsor.
Because $B_T$ is free over $T[[x,y]]$, the sequence $x,y$ is
regular on $B_T$.  The two-open \v Cech calculation gives
\[
 \Gamma(V_T,\mathcal O)=B_T.
\]
In particular, clopen decompositions of $V_T$ are the idempotent
decompositions of $B_T$.  Completeness lifts idempotents uniquely
from $B_T/(x,y)B_T$; the latter is
$(B/(x,y)B)\otimes_kT$.  Since $B/(x,y)B$ is a local Artin
$k$-algebra with residue field $k$, its nilpotent augmentation
ideal shows that these idempotents are exactly those of $T$
\cite[Section 15.11, tags 09XD, 09XI and 0ALI]{C_StacksFormalLift}.
Every automorphism of the finite \'etale cover $V_T/U_T$
is given by a locally constant deck element on $V_T$; by the
preceding idempotent calculation that element is locally constant
on $\operatorname{Spec}T$.  Thus it belongs to $G(T)$.
The Hartogs equality makes restriction injective, and all such
deck elements extend to $B_T$.  This proves the claim on every
test algebra, not just on geometric points.

The ring $A$ is excellent and strictly henselian.  Gabber's local
finiteness theorem says that $H^1(U,G)$ is finite for every finite
group $G$, where $U$ is its punctured spectrum
\cite[XXI, Theorem 1.4, p. 529]{ILOXXID}.
The divisible abelian group $k^*$ acts on this finite set by
scaling.  Its image in a finite permutation group is both
divisible and finite, hence trivial.  Each scaling therefore
lifts to the cover and to its normal finite algebra $B$.
It follows that $E(k)\to k^*$ is surjective.  In particular,
$k[t,t^{-1}]\to\mathcal O(E)$ is injective.

Faithful flatness of a commutative Hopf algebra over any Hopf
subalgebra, with no finite-type restriction, now makes $\rho$
faithfully flat
\cite[VI, Theorem 11.1, pp. 82--85]{C_MilneAGS2012}.
It is affine, hence fpqc.  The identity
$E\times_{\mathbb G_m}E\simeq E\times G$ and fpqc descent
show that $\rho$ is finite \'etale.  Its identity component
$E^0$ is a connected one-dimensional affine smooth group with
a finite-kernel surjection to $\mathbb G_m$; it is a torus.
Indeed, the only connected smooth affine groups of dimension
one over $k$ are $\mathbb G_a$ and $\mathbb G_m$, and
$\mathbb G_a$ has no nontrivial character
\cite[Corollary 16.16, p. 329]{C_MilneAG2022}.
After choosing its coordinate, $E^0\to\mathbb G_m$ is
$t\mapsto t^n$ with $n>0$ prime to $p$.  Conjugation by this
connected group on the constant finite group $G$ is trivial.
This proves the required lift and centralization.
\end{proof}

\begin{proposition}[Proved: radial tameness for completed graded surface cones]
\label{D:graded-cone-tame}
In the situation of Lemma~\ref{D:formal-torus-lift}, assume the
positive grading is effective, so that its occurring degrees
have greatest common divisor one.  Every extension to
$\operatorname{Frac}(B)$ of the radial divisorial valuation of
$\operatorname{Frac}(A)$ is tame.  Thus a normal surface germ
whose completed resolution along its exceptional curve is the
completion of the zero section in a negative line bundle over
a smooth proper curve has no wild inertia at that exceptional
valuation in any finite cover \'etale on the punctured germ.
\end{proposition}
\begin{proof}
Let $\widetilde T=E^0$ be the lifted torus, acting on $A$ with
weights $n$ times the original weights.  The finite-dimensional
quotients $B/\mathfrak m_A^rB$ are rational torus
representations.  Exactness of weight spaces lets us lift a
finite weight basis of $B/\mathfrak m_AB$ compatibly through
all these quotients.  We obtain eigenvectors $b_1,\ldots,b_s$
which generate the finite $A$-module $B$ by Nakayama's lemma.
If $w_j$ is the weight of $b_j$, projection of an expression
$b=\sum a_jb_j$ to a weight $d$ shows that every eigenvector
of weight $d$ is an $R$-linear combination of these $b_j$.
The coefficient of $b_j$ is the homogeneous component of
$a_j$ with torus weight $d-w_j$, which lies in $R$.

It follows that the direct sum $S$ of the weight spaces in $B$
is a finite graded $R$-algebra and that its weights are bounded
below.  A nonzero element of negative weight would have powers
of arbitrarily negative weight, contradicting this bound and
the fact that $B$ is a domain.  Thus $S$ is nonnegatively graded.
Its degree-zero part is a finite-dimensional domain over the
algebraically closed field $k$, hence equals $k$.
The same eigenvector expansions identify $B$ with the product
of its nonnegative weight spaces: the weight filtration and
the maximal-ideal filtration are cofinal because $S$ is finite
over the positively graded ring $R$.  In particular, $B$ is
the completion of $S$ at its positive ideal.

The group $G$ commutes with $\widetilde T$, hence preserves $S$
and every weight space.  Since $B^G=A$, taking weight spaces
gives $S^G=R$.  Put $L_0=\operatorname{Frac}(S)$ and
$K_0=\operatorname{Frac}(R)$.  Then $L_0/K_0$ is Galois with
group $G$: the action is faithful because $S$ is dense in $B$,
and invariants of the fraction field are fractions of invariants.

Let $g$ be the greatest common divisor of the positive weights
occurring in $S$.  Lowest weight divided by $g$ defines the
radial discrete valuation on $L_0$, and also on
$\operatorname{Frac}(B)$ by the completed weight expansion.
It is preserved by $G$.  Its restriction to the base radial
valuation has ramification index
\[
 e=n/g.
\]
Indeed, the weights of $R$ in $S$ are $n$ times its effective
original weights.  Thus $g\mid n$, and $e$ is prime to $p$.
The residue field of the radial valuation on $L_0$, or on the
completed fraction field, is the field of fractions of
homogeneous expressions of equal degree; equivalently it is
$L_0^{\widetilde T}$.  On the base it is
$K_0^{\mathbb G_m}$.  Since the two actions commute,
\[
 \bigl(L_0^{\widetilde T}\bigr)^G
   =\bigl(L_0^G\bigr)^{\widetilde T}
   =K_0^{\mathbb G_m}.
\]
The residue extension is therefore separable, being an extension
of a fixed field under a finite group.  The exhibited valuation
is $G$-invariant; transitivity of $G$ on prolongations in the
Galois extension shows it is the only prolongation.  Its
prime-to-$p$ ramification index and separable residue extension
prove tameness.

For the final assertion, a negative line bundle $M$ on a smooth
proper curve $C$ has section ring
$R=\bigoplus_{d\geq0}H^0(C,M^{-d})$.  This is a finitely
generated normal ring.  Finite generation follows from ampleness
of $M^{-1}$; normality follows by intersecting the corresponding
divisorial valuation rings inside $k(C)(t)$.  The functions on
the formal zero-section neighborhood are
$\prod_{d\geq0}H^0(C,M^{-d})$, which is the completion of $R$
at its positive ideal.  For a proper normal surface contraction,
degree-zero formal functions identifies this ring with the
completed local ring downstairs: the exceptional-ideal and
maximal-ideal topologies on the resolution are cofinal
\cite[Theorem 30.20.5, tag 02OC]{C_StacksFormalPic}.
Its radial valuation is exactly the valuation of $C$.
The degrees are effective since all sufficiently large degrees
have nonzero sections.  A finite cover of the original
punctured germ pulls back to a disjoint union of connected
finite covers of this completed puncture.  Each is dominated
by a Galois cover, to which the preceding argument applies.
The formal line-bundle identification supplies a common genuine
completed boundary DVR $k(C)[[t]]$: both the original boundary
DVR and the radial DVR on the vertex completion have this
completion, compatibly with the maps from their function fields.
We do not identify vertex completion with DVR completion.
Completion preserves the ramification index and separability
of the residue extension in a finite extension of excellent
DVRs.  Testing the cover over this common complete DVR therefore
descends tameness to the original boundary valuation.
The same argument applies when
we identify the completed resolved germ with this formal
line-bundle neighborhood.
\end{proof}
